\documentclass[10pt,a4paper]{amsart}
\usepackage{amsmath,amssymb,amsthm,graphicx}
\usepackage{amsmath}
\usepackage{amssymb}
\usepackage{hyperref}
\usepackage[all]{xy}
\numberwithin{equation}{section}
\newtheorem*{theointro}{Theorem}
\newtheorem{Theorem}{Theorem}[section]
\newtheorem{Corollary}[Theorem]{Corollary}
\newtheorem{Lemma}[Theorem]{Lemma}
\newtheorem{Proposition}[Theorem]{Proposition}
{ \theoremstyle{definition}
\newtheorem{Definition}[Theorem]{Definition}
\newtheorem{Remark}[Theorem]{Remark} }
\newcommand{\divi}[1]{\operatorname{div}_{#1}}
\newcommand{\weight}[1]{\omega_{#1}}
\newcommand{\degr}[1]{\deg_{#1}}
\newcommand{\wpf}[1]{\wp_{#1}}
\newcommand{\prodtheta}[1]{\Theta_{#1}}
\newcommand{\quotprodtheta}[1]{\Theta_{#1}^{\rm quot}}
\title{A bound on the degrees of the elliptic decomposition of a function with N poles}
\author{Thomas Dreyfus, Julien Roques}
\date{doi:10.3842/SIGMA.2015.003; CC BY4.0}
\begin{document}
\maketitle
\noindent\textit{Editorial note.} All delivered mathematical text is the author's own, copied from the pinned publisher source sources/sigma-2015-003-elliptic-difference-galois/pinned-original-source.tex; no statement or proof step was rewritten, reordered or removed.\par\medskip
\noindent\textit{Editorial note.} Preamble provenance: the retained preamble is the source's own layout adapter, if any, plus the author's macro block at lines 9-18 of sources/sigma-2015-003-elliptic-difference-galois/pinned-original-source.tex (everything between the publisher class metadata and \textbackslash{}begin\{document\}, comments excluded). Nothing was added, no macro was invented, and the publisher document class is deliberately not reproduced; the excerpt serves only as preamble, never as a span coordinate.\par\medskip
\noindent\textit{Editorial note.} Metadata provenance: source\_id and source\_version were taken from sources/sigma-2015-003-elliptic-difference-galois/license-observation.json; the title from SOURCE.json; the author list from the source metadata record; the source URL from the first non-empty of source\_url, source\_urls[0] and the evidence's fixed\_version\_url/license\_url. The licence evidence file used by this card is sources/sigma-2015-003-elliptic-difference-galois/license-observation.json, recorded in license.evidence\_path so that the binding target is explicit.\par\medskip
\noindent\textit{Editorial note.} External references are navigation pointers only; every deliberately kept step of the author's argument stays inside the counted spans. The displayed numbers of those pointers are the author's own numbering under the author's \textbackslash{}newtheorem declarations (section.index as declared), computed mechanically by the lane's scanner.\par\medskip
\noindent\textit{Editorial note.} The counted proof contains no reference to any labelled item and no bibliography citation, so the card needs neither a navigation stub nor a bibliography excerpt.\par\medskip
\noindent\textit{Editorial note.} This is the first card this lane authors from this source, taken from the newly approved batch and selected as a zero-stub, zero-citation candidate to keep the publishable queue fed.\par\medskip
\noindent\textit{Editorial note.} Assembly history: r001 failed with 'Undefined control sequence' at \textbackslash{}prodtheta because this source declares its math macros (lines 22-27) AFTER \textbackslash{}begin\{document\}, so a preamble span that stops at \textbackslash{}begin\{document\} misses them. r002 adds that run as a second preamble span (pure command definitions, copied verbatim); the generator now detects such post-document macro runs automatically. No mathematics changed.\par\medskip
\noindent\textit{Allowed standard prerequisites (editorial).}\begin{itemize}
\item The identification of functions periodic with respect to k Lambda with Lambda-periodic functions by scaling, and the even/odd decomposition formulas for f in terms of A(weierstrass) and weierstrass' times B. Used inside the counted proof at lines 627-630 and 631-634, 643-646. Reduces to k = 1 and produces the explicit expressions whose pole counts are estimated.
\item Elliptic functions have as many zeroes as poles, counted with multiplicity. Used inside the counted proof at lines 640-641. Turns the pole bound on A(weierstrass) into the degree bound on its numerator.
\end{itemize}\medskip
\begin{Lemma}
\label{lem2}
Assume that $f \in \operatorname{M}_{k\Lambda}$, seen has a~meromorphic function over $\mathbb C/k\Lambda$, has at
most~$N$ poles counted with multiplicities $($or, equivalently, that $f=p/q$ with $p,q \in \prodtheta{k}$ such that
${\degr{k} p, \degr{k} q \leq N})$.
Then, there exist $A=P/Q$ and $B=R/S$ with~$P,Q \in \mathbb C[X]$ of degree at most $2N$ and $R,S \in \mathbb C[X]$ of
degree at most~$2N+3$ such that
\begin{gather*}
f=A(\wpf{k})+\wpf{k}' B(\wpf{k}).
\end{gather*}
\end{Lemma}

\begin{proof}
Using the fact that $f(z)$ belongs to $\operatorname{M}_{k\Lambda}$ if and only if $f(kz)$ belongs to
$\operatorname{M}_{\Lambda}$, it is easily seen that it is suf\/f\/icient to prove the lemma for $k=1$.
In what follows, we see the~$\Lambda$-periodic elliptic functions as meromorphic functions on $\mathbb C/\Lambda$.
Let~$A,B \in \mathbb C(X)$ be such that $f=A(\wp)+\wp' B(\wp)$.
It follows from the formula
\begin{gather*}
A(\wp(z))=\frac{f(z)+f(-z)}{2}
\end{gather*}
that $A(\wp)$ has at most $2N$ poles counted with multiplicities in $\mathbb C/\Lambda$.
But, if $A=P/Q$ with $\gcd(P,Q)=1$ then $A(\wp)$ has at least $\deg Q$ poles counted with multiplicities in $\mathbb
C/\Lambda$ (namely, the zeroes of $Q(\wp)$).
So $\deg Q \leq 2N$.

Using the fact that elliptic functions have the same numbers of zeroes and poles, the same argument applied to
$1/A(\wp)$ shows that~${\deg P \leq 2N}$.

Using the formula
\begin{gather*}
B(\wp(z))=\frac{f(z)-f(-z)}{2 \wp'(z)},
\end{gather*}
similar arguments show that $\deg R \leq 2N+3$ and $\deg S \leq 2N+3$.
\end{proof}

\end{document}
